%% toc: P3 41 41.2
\section{Proofs for nilpotent groups of class two and further group examples}
\label{sec:annex-p03-w03-p03-074-nilpotent}

\begin{annexitem}{annex:p03-w03-p03-074-class-law}{thm:p03-w03-p03-074-class-law}{Each nilpotency class is axiomatized by one universal law}
\begin{proof}
\emph{Item 1.} Suppose \(Z_c(G)=G\). We show by downward induction on
\(j=c,c-1,\ldots,0\) that every left-normed commutator
\([g_1,\ldots,g_{c+1-j}]\) of elements of \(G\) lies in \(Z_j(G)\). For
\(j=c\) the commutator is the single element \(g_1\), which lies in
\(Z_c(G)=G\). If the claim holds for some \(j\ge1\), let
\(u=[g_1,\ldots,g_{c+1-j}]\in Z_j(G)\) and \(h\in G\). By the definition
of \(Z_j(G)\), \([u,h]\in Z_{j-1}(G)\), and
\([u,h]=[g_1,\ldots,g_{c+1-j},h]\) is an arbitrary commutator of weight
\(c+2-j=c+1-(j-1)\). At \(j=0\) every commutator of weight \(c+1\) lies in
\(Z_0(G)=\{1\}\), which is \(G\models\lambda_{c+1}\).

Conversely suppose \(G\models\lambda_{c+1}\). We show by upward induction
on \(j=0,1,\ldots,c\) that every commutator of weight \(c+1-j\) lies in
\(Z_j(G)\). For \(j=0\) this is the hypothesis. If it holds for some
\(j<c\), let \(u=[g_1,\ldots,g_{c-j}]\) be a commutator of weight
\(c+1-(j+1)\). For every \(h\in G\), \([u,h]\) is a commutator of weight
\(c+1-j\) and so lies in \(Z_j(G)\); by definition \(u\in Z_{j+1}(G)\). At
\(j=c\) every commutator of weight \(1\), that is every element of \(G\),
lies in \(Z_c(G)\), so \(Z_c(G)=G\).

\emph{Item 2.} If \(G\models\lambda_{c+1}\), then for all
\(g_1,\ldots,g_{c+2}\),
\[
 [g_1,\ldots,g_{c+2}]=[[g_1,\ldots,g_{c+1}],g_{c+2}]=[1,g_{c+2}]
 =1\cdot g_{c+2}^{-1}\cdot1\cdot g_{c+2}=1.
\] By item 1, \(Z_c(G)=G\)
implies \(Z_{c+1}(G)=G\), and induction gives \(Z_{c'}(G)=G\) for every
\(c'\ge c\).

\emph{Item 3.} The first sentence is item 1. Let \(G\models\lambda_{c+1}\).
A subgroup \(H\) is a substructure for \(\{\cdot,{}^{-1},1\}\), and a
commutator term evaluated at elements of \(H\) has the same value in \(H\)
as in \(G\); so \(H\models\lambda_{c+1}\), and \(H\) satisfies the group
axioms, which are universal. If \(\rho:G\to H\) is a surjective
homomorphism, every tuple \(h_1,\ldots,h_{c+1}\) of \(H\) has the form
\(\rho g_1,\ldots,\rho g_{c+1}\), and since \(\rho\) preserves products and
inverses,
\([h_1,\ldots,h_{c+1}]=\rho[g_1,\ldots,g_{c+1}]=\rho(1)=1\). In a direct
product the operations are coordinatewise, so a commutator of tuples is the
tuple of the coordinate commutators, and it is the identity when every
coordinate is. For an ultraproduct \(\prod_iG_i/\mathcal U\) of groups
satisfying \(\lambda_{c+1}\), each group axiom and \(\lambda_{c+1}\) hold
in every factor, the index set belongs to \(\mathcal U\), and
\cref{thm:los} gives the same sentences in the ultraproduct.
\end{proof}
\end{annexitem}

\begin{annexitem}{annex:p03-w03-p03-074-centre-step}{thm:p03-w03-p03-074-centre-step}{Passing to the central quotient lowers the weight by one}
\begin{proof}
\emph{Item 1.} \(G\models\lambda_{k+1}\) says that
\([[g_1,\ldots,g_k],h]=1\) for all \(g_1,\ldots,g_k,h\). Since
\([u,h]=1\) exactly when \(uh=hu\), this says that each
\([g_1,\ldots,g_k]\) commutes with every \(h\), that is, lies in
\(Z(G)\).

\emph{Item 2.} A homomorphism preserves products and inverses, hence
\[
 \rho[g_1,\ldots,g_k]=[\rho g_1,\ldots,\rho g_k].
\] Because \(\rho\) is
surjective, every \(k\)-tuple of \(H\) is the image of a \(k\)-tuple of
\(G\), obtained by choosing one preimage of each of its finitely many
entries. Hence \(H\models\lambda_k\) if and only if
\(\rho[g_1,\ldots,g_k]=1\) for all \(g_1,\ldots,g_k\in G\), if and only if
\([g_1,\ldots,g_k]\in\ker\rho=Z(G)\) for all such tuples, if and only if
\(G\models\lambda_{k+1}\) by item 1.
\end{proof}
\end{annexitem}

\begin{annexitem}{annex:p03-w03-p03-074-class-two-rules}{thm:p03-w03-p03-074-class-two-rules}{Computing rules in class two}
\begin{proof}
The three general identities are direct computations:
\[\begin{aligned}
 [x,y]^{-1}&=(x^{-1}y^{-1}xy)^{-1}=y^{-1}x^{-1}yx=[y,x],\\
 xy[y,x]&=xy\,y^{-1}x^{-1}yx=yx,\\
 {}[x,z]^y[y,z]&=y^{-1}x^{-1}z^{-1}xzy\,y^{-1}z^{-1}yz
   =y^{-1}x^{-1}z^{-1}xyz=[xy,z].
\end{aligned}\]

Now let \(G\models\lambda_3\).

\emph{Item 1.} \(\lambda_3\) says \([[x,y],z]=1\) for every \(z\), so
\([x,y]\) commutes with every element.

\emph{Item 2.} By item 1, \([x,z]\) is central, so \([x,z]^y=[x,z]\) and
the third general identity gives \([xy,z]=[x,z][y,z]\). Using the first
general identity, this rule, and the fact that central elements commute,
\[\begin{aligned}
 [x,yz]&=[yz,x]^{-1}=([y,x][z,x])^{-1}\\
 &=[z,x]^{-1}[y,x]^{-1}=[x,z][x,y]=[x,y][x,z].
\end{aligned}\]

\emph{Item 3.} For \(m\ge0\) argue by induction: \([x^0,y]=[1,y]=1\), and
\([x^{m+1},y]=[x^mx,y]=[x^m,y][x,y]=[x,y]^{m+1}\) by item 2. Next,
\([x^{-1},y][x,y]=[x^{-1}x,y]=1\), so \([x^{-1},y]=[x,y]^{-1}\), and for
\(m\ge0\),
\[
 [x^{-m},y]=[(x^{-1})^m,y]=[x^{-1},y]^m=[x,y]^{-m}.
\]
Finally
\[
 [x,y^m]=[y^m,x]^{-1}=([y,x]^m)^{-1}=(([x,y]^{-1})^m)^{-1}=[x,y]^m.
\]

\emph{Item 4.} For \(n=0\) both sides are \(1\). Suppose the identity holds
for \(n\) and put \(N=n(n-1)/2\). The second general identity with \(y^n\)
in place of \(y\) gives \(y^nx=xy^n[y^n,x]\), and \([y^n,x]=[y,x]^n\) by
item 3. Since \([y,x]\) is central,
\[\begin{aligned}
 (xy)^{n+1}&=x^ny^n[y,x]^Nxy
 =x^n(y^nx)y[y,x]^N\\
 &=x^nxy^n[y,x]^ny[y,x]^N
 =x^{n+1}y^{n+1}[y,x]^{n+N},
\end{aligned}\]
and \(n+N=n(n+1)/2\).
\end{proof}
\end{annexitem}

\begin{annexitem}{annex:p03-w03-p03-074-heisenberg-basic}{thm:p03-w03-p03-074-heisenberg-basic}{Structure of the Heisenberg group}
\begin{proof}
\emph{Item 1.} For \(g=(a,b,c)\), \(h=(a',b',c')\) and
\(k=(a'',b'',c'')\), both \((gh)k\) and \(g(hk)\) have first coordinate
\(a+a'+a''\), second coordinate \(b+b'+b''\), and third coordinate
\(c+c'+c''+ab'+ab''+a'b''\): in \((gh)k\) it is
\((c+c'+ab')+c''+(a+a')b''\), and in \(g(hk)\) it is
\(c+(c'+c''+a'b'')+a(b'+b'')\). The element \((0,0,0)\) is a two-sided
identity since its first two coordinates are \(0\). Further
\((a,b,c)(-a,-b,-c+ab)=(0,0,c-c+ab-ab)\) and
\((-a,-b,-c+ab)(a,b,c)=(0,0,-c+ab+c-ab)\), both the identity. The
first two coordinates of a product are sums, so \(\pi\) is a homomorphism,
and \(\pi(a,b,0)=(a,b)\) shows that it is surjective. Finally
\(\zeta(c)\zeta(c')=(0,0,c+c'+0\cdot0)=\zeta(c+c')\), and \(\zeta\) is
injective.

\emph{Item 2.} For every \(e\in R\) and \((u,v,w)\in\mathrm H(R)\),
\((u,v,w)\zeta(e)=(u,v,w+e+u\cdot0)\) and
\(\zeta(e)(u,v,w)=(u,v,e+w+0\cdot v)\), so \(\zeta(e)\) is central. With
\(g,h\) as above, \(gh=(a+a',b+b',c+c'+ab')\) and
\(hg=(a+a',b+b',c+c'+a'b)\), so \(gh=hg\,\zeta(ab'-a'b)\). Hence
\([g,h]=g^{-1}h^{-1}gh=(hg)^{-1}gh=\zeta(\Delta(g,h))\).

\emph{Item 3.} Item 2 shows \(\zeta(R)\subseteq Z(\mathrm H(R))\), and
\(\zeta(R)=\ker\pi\) by the definitions. If \(g=(a,b,c)\) is central,
then \(1=[g,e_2]=\zeta(a\cdot1-0\cdot b)=\zeta(a)\) and
\(1=[g,e_1]=\zeta(a\cdot0-1\cdot b)=\zeta(-b)\); by injectivity of
\(\zeta\), \(a=b=0\) and \(g\in\zeta(R)\). For the centralizer,
\(gh=hg\) iff \([g,h]=1\) iff \(\zeta(\Delta(g,h))=\zeta(0)\) iff
\(\Delta(g,h)=0\).

\emph{Item 4.} Every commutator lies in \(\zeta(R)=Z(\mathrm H(R))\) by
items 2 and 3, so \(\mathrm H(R)\models\lambda_3\) by
\cref{thm:p03-w03-p03-074-centre-step}, item 1, with \(k=2\). If
\(1\neq0\) in \(R\), then \([e_1,e_2]=\zeta(1\cdot1-0\cdot0)=\zeta(1)\)
is not the identity, so \(\lambda_2\) fails, and \(\lambda_1\) fails since
the group is nontrivial. By \cref{thm:p03-w03-p03-074-class-law}, item 1,
\(Z_0\) and \(Z_1\) are proper and \(Z_2\) is the whole group, so the
class is exactly \(2\).

\emph{Item 5.} For \(n=0\) both sides are \((0,0,0)\). If the formula holds
for \(n\), then
\((a,b,c)^{n+1}=(na,nb,nc+\tfrac{n(n-1)}2ab)(a,b,c)
=((n+1)a,(n+1)b,(n+1)c+(\tfrac{n(n-1)}2+n)ab)\), using the product rule
with \(na\cdot b=n\,ab\); and \(\tfrac{n(n-1)}2+n=\tfrac{(n+1)n}2\).
\end{proof}
\end{annexitem}

\begin{annexitem}{annex:p03-w03-p03-074-ring-recovery}{thm:p03-w03-p03-074-ring-recovery}{The ring inside its Heisenberg group}
\begin{proof}
Write \(\det(A,B)=A_1B_2-A_2B_1\) for \(A,B\in R^2\). In a commutative
ring \(\det\) is linear in each argument, \(\det(A,A)=0\) and
\(\det(B,A)=-\det(A,B)\); and \(\Delta(g,h)=\det(\pi g,\pi h)\).

\emph{Step 1 (Cramer identity).} Let \(A,B\in R^2\) with
\(d=\det(A,B)\) a unit. Then every \(w\in R^2\) satisfies
\[
 w=\det(w,B)d^{-1}A+\det(A,w)d^{-1}B .
\]
Let \(w'\) be \(w\) minus the right side. By linearity,
\(\det(w',B)=\det(w,B)-\det(w,B)d^{-1}d-0=0\) and
\(\det(A,w')=\det(A,w)-0-\det(A,w)d^{-1}d=0\). Thus
\(w'_1B_2=w'_2B_1\) and \(A_1w'_2=A_2w'_1\). Then
\[
 d\,w'_1=A_1B_2w'_1-A_2B_1w'_1=A_1B_1w'_2-B_1A_1w'_2=0,
\]
using the first equation in the first term and the second in the second;
and
\[
 d\,w'_2=A_1B_2w'_2-A_2B_1w'_2=B_2A_2w'_1-A_2B_2w'_1=0,
\]
using the second equation in the first term and the first in the second.
As \(d\) is a unit, \(w'=0\).

\emph{Step 2 (centralizers).} Put \(X=\pi x\), \(Y=\pi y\), so
\(\delta=\det(X,Y)\). For \(s\in\mathrm H(R)\),
\([s,x]=1\) iff \(\det(\pi s,X)=0\) by
\cref{thm:p03-w03-p03-074-heisenberg-basic}, items 2 and 3. If
\(\pi s=\alpha X\) then \(\det(\alpha X,X)=0\). Conversely, if
\(\det(\pi s,X)=0\) then \(\det(X,\pi s)=0\), and Step 1 with
\(A=X\), \(B=Y\) gives \(\pi s=\det(\pi s,Y)\delta^{-1}X\). So
\([s,x]=1\) iff \(\pi s\in RX\). In the same way, \([t,y]=1\) iff
\(\det(\pi t,Y)=0\) iff \(\pi t=\det(X,\pi t)\delta^{-1}Y\in RY\).

\emph{Step 3 (the multiplication relation).} Let
\(r_1,r_2,r_3\in R\). Suppose \(s,t\) witness
\(\mathrm{Mul}(x,y;\zeta(r_1),\zeta(r_2),\zeta(r_3))\). By Step 2,
\(\pi s=\alpha X\) and \(\pi t=\beta Y\) for some \(\alpha,\beta\). Then
\(\zeta(r_1)=[s,y]=\zeta(\det(\alpha X,Y))=\zeta(\alpha\delta)\), so
\(\alpha=r_1\delta^{-1}\); likewise
\(\zeta(r_2)=[x,t]=\zeta(\beta\delta)\), so \(\beta=r_2\delta^{-1}\); and
\(\zeta(r_3)=[s,t]=\zeta(\alpha\beta\delta)\), so
\(r_3=r_1r_2\delta^{-1}\). Conversely, for \(r_3=r_1r_2\delta^{-1}\), put
\(\alpha=r_1\delta^{-1}\), \(\beta=r_2\delta^{-1}\),
\(s=(\alpha X_1,\alpha X_2,0)\) and \(t=(\beta Y_1,\beta Y_2,0)\). Then
\([s,x]=\zeta(\alpha\det(X,X))=1\), \([t,y]=1\), \([s,y]=\zeta(r_1)\),
\([x,t]=\zeta(r_2)\) and \([s,t]=\zeta(r_3)\) by the same computations.
Since every central element has the form \(\zeta(r)\)
(\cref{thm:p03-w03-p03-074-heisenberg-basic}, item 3), for central
\(z_1,z_2\) there is exactly one \(z_3\) with
\(\mathrm{Mul}(x,y;z_1,z_2,z_3)\), namely
\(\zeta(r_1r_2\delta^{-1})\) when \(z_i=\zeta(r_i)\), and it is central.

\emph{Step 4 (interpretation).} The conditions of
\cref{def:p03-imaginaries-interpretation} are statements about
\((\mathrm H(R),x,y)\), so they are consequences of its complete theory
as soon as they are true in it. The domain \(\mathrm{Ctr}\) defines
\(Z(\mathrm H(R))=\zeta(R)\), which contains the identity. Equality is
literal, so it is an equivalence relation and every graph is invariant.
The graph \(z=1\) has the unique central output \(\zeta(0)\); the graph
\(z=[c_x,c_y]\) has the unique output \(\zeta(\delta)\), central by
\cref{thm:p03-w03-p03-074-heisenberg-basic}, item 2; the graph
\(z_3=z_1z_2\) has the unique output \(\zeta(r_1+r_2)\) on inputs
\(\zeta(r_1),\zeta(r_2)\), central; and the graph of \(\cdot\) is total and
single-valued with central outputs by Step 3. No target theory is imposed.
So \(I\) is an interpretation, and \(I(\mathrm H(R),x,y)\) is the
\(\Lang_{\mathrm{rg}}\)-structure on \(\zeta(R)\) with \(0=\zeta(0)\),
\(1=\zeta(\delta)\), \(\zeta(r_1)+\zeta(r_2)=\zeta(r_1+r_2)\) and
\(\zeta(r_1)\cdot\zeta(r_2)=\zeta(r_1r_2\delta^{-1})\).

\emph{Step 5 (isomorphism).} The map \(\theta_{x,y}\) is a bijection
\(\zeta(R)\to R\), because \(\zeta\) is a bijection onto \(\zeta(R)\) and
multiplication by the unit \(\delta^{-1}\) is a bijection of \(R\). It
sends \(\zeta(0)\) to \(0\) and \(\zeta(\delta)\) to \(1\); it sends
\(\zeta(r_1+r_2)\) to \(r_1\delta^{-1}+r_2\delta^{-1}\); and it sends
\(\zeta(r_1r_2\delta^{-1})\) to
\(r_1r_2\delta^{-2}=(r_1\delta^{-1})(r_2\delta^{-1})\). So it is an
isomorphism.

\emph{Step 6 (sentences).} Let \(\sigma\) be an
\(\Lang_{\mathrm{rg}}\)-sentence. An isomorphism preserves the truth of
every formula at corresponding assignments, by induction on formulas: the
atomic case holds because the isomorphism commutes with the constants and
operations and is injective, the Boolean cases are immediate, and the
quantifier case uses that it is a bijection. Hence \(R\models\sigma\) iff
\(I(\mathrm H(R),x,y)\models\sigma\). By
\cref{thm:p03-imaginaries-quotient-truth} with the empty assignment, this
holds iff \((\mathrm H(R),x,y)\models\sigma^I\), where \(\sigma^I\)
contains the constants \(c_x,c_y\). The constants denote \(x,y\), so this
holds iff \(\mathrm H(R)\models\sigma^I(x,y)\) with the fresh variables
assigned \(x\) and \(y\).
\end{proof}
\end{annexitem}

\begin{annexitem}{annex:p03-w03-p03-074-unit-criterion}{thm:p03-w03-p03-074-unit-criterion}{The good parameter pairs are definable}
\begin{proof}
Use the notation of the annex proof of
\cref{thm:p03-w03-p03-074-ring-recovery}, with \(X=\pi x\), \(Y=\pi y\)
and \(\delta=\det(X,Y)=\Delta(x,y)\).

Suppose \(\delta\) is a unit. A central \(z\) has the form \(\zeta(r)\).
Put
\[
 s=(r\delta^{-1}X_1,\,r\delta^{-1}X_2,\,0),\qquad
 t=(r\delta^{-1}Y_1,\,r\delta^{-1}Y_2,\,0).
\]
By \cref{thm:p03-w03-p03-074-heisenberg-basic}, item 2,
\[
 [s,x]=\zeta(r\delta^{-1}\det(X,X))=1,\qquad
 [s,y]=\zeta(r\delta^{-1}\det(X,Y))=\zeta(r)=z,
\]
\[
 [t,y]=\zeta(r\delta^{-1}\det(Y,Y))=1,\qquad
 [x,t]=\zeta(r\delta^{-1}\det(X,Y))=z.
\]
Both conjuncts of \(\beta(x,y)\) hold.

Conversely suppose \(\mathrm H(R)\models\beta(x,y)\). Apply both
conjuncts to the central element \(\zeta(1)\). This gives \(s\) with
\(\det(\pi s,X)=0\) and \(\det(\pi s,Y)=1\), and \(t\) with
\(\det(\pi t,Y)=0\) and \(\det(X,\pi t)=1\). Put \(S=\pi s\) and
\(T=\pi t\). Since \(\det(S,Y)=1\) is a unit, Step 1 of the annex proof
of \cref{thm:p03-w03-p03-074-ring-recovery}, with \(A=S\), \(B=Y\) and
\(w=X\), gives
\[
 X=\det(X,Y)S+\det(S,X)Y=\delta S .
\]
Then \(1=\det(X,T)=\det(\delta S,T)=\delta\det(S,T)\), so \(\delta\) is a
unit with inverse \(\det(S,T)\).

For the last sentence, \(\Delta(e_1,e_2)=1\cdot1-0\cdot0=1\).
\end{proof}
\end{annexitem}

\begin{annexitem}{annex:p03-w03-p03-074-equivalence}{thm:p03-w03-p03-074-equivalence}{A commutative ring and its Heisenberg group have equivalent theories}
\begin{proof}
\emph{Item 1.} Suppose \(R\models\sigma\). By
\cref{thm:p03-w03-p03-074-unit-criterion},
\(\mathrm H(R)\models\beta(e_1,e_2)\), and since \(\Delta(e_1,e_2)=1\),
\cref{thm:p03-w03-p03-074-ring-recovery} gives
\(\mathrm H(R)\models\sigma^I(e_1,e_2)\). So \(e_1,e_2\) witness
\(\sigma^{\sharp}\). Conversely, suppose \(x,y\) witness
\(\sigma^{\sharp}\). By \cref{thm:p03-w03-p03-074-unit-criterion},
\(\Delta(x,y)\) is a unit, and \cref{thm:p03-w03-p03-074-ring-recovery}
turns \(\mathrm H(R)\models\sigma^I(x,y)\) into \(R\models\sigma\).
The formula \(\beta\) is fixed, and \(\sigma^I\) is computed by the
recursion of \cref{def:p03-imaginaries-translation} from the five fixed
formulas assigned by \(I\), with bound variables renamed by a fixed
computable rule; none of these depends on \(R\).

\emph{Item 2.} Every triple satisfies the domain formula, and
coordinatewise equality is an equivalence relation whose classes are
singletons, so the quotient map identifies \(R^3\) with the interpreted
domain and every graph is invariant. The identity graph has the unique
output \((0,0,0)\). For input \((a,b,c)\) the inverse graph has the unique
output \((-a,-b,ab-c)\), which is \((a,b,c)^{-1}\) by
\cref{thm:p03-w03-p03-074-heisenberg-basic}, item 1. For inputs
\((a,b,c),(a',b',c')\) the product graph has the unique output
\((a+a',b+b',c+c'+ab')\), the product in \(\mathrm H(R)\). These facts are
true in \(R\) and therefore follow from \(\operatorname{Th}(R)\). The
identity of \(R^3\) is thus an isomorphism of the interpreted structure
with \(\mathrm H(R)\), and the interpretation is computable because both
signatures are finite.

\emph{Item 3.} Item 1 is a computable many-one reduction of
\(\operatorname{Th}(R)\) to \(\operatorname{Th}(\mathrm H(R))\). Item 2 and
\cref{thm:p03-w03-p03-022-transfer} give the reduction
\(\tau\mapsto\tau^J\) of \(\operatorname{Th}(\mathrm H(R))\) to
\(\operatorname{Th}(R)\), where \(J\) is the interpretation of item 2;
its formulas do not depend on \(R\). If one theory is decidable, composing
its decision procedure with the reduction into it decides the other.
\end{proof}
\end{annexitem}

\begin{annexitem}{annex:p03-w03-p03-074-examples}{thm:p03-w03-p03-074-examples}{Decidable and undecidable Heisenberg groups}
\begin{proof}
\emph{Item 1.} The theory of \((\mathbb Z,+,\cdot,0,1)\), whose signature
is \(\Lang_{\mathrm{rg}}\), is undecidable by
\cref{thm:p02-w02c-p02-003-integers}. By
\cref{thm:p03-w03-p03-074-equivalence}, item 3, the theory of
\(\mathrm H(\mathbb Z)\) is undecidable.

\emph{Item 2.} A field is a commutative ring with identity, so
\cref{thm:p03-w03-p03-074-equivalence}, item 3, applies. If \(K\) is a
finite field, so is the structure \(\mathrm H(K)\), with \(|K|^3\)
elements and computable operation tables. The truth of a sentence in a
finite structure is computed by recursion on the sentence: atomic
sentences with names of elements are evaluated from the tables, Boolean
connectives by their truth tables, and a quantifier by evaluating its
matrix at each of the finitely many elements. Hence the theory of
\(\mathrm H(K)\) is decidable; the same recursion decides the theory of
\(K\), in agreement with item 3 of
\cref{thm:p03-w03-p03-074-equivalence}.
\end{proof}
\end{annexitem}

\begin{annexitem}{annex:p03-w03-p03-074-dihedral-class}{thm:p03-w03-p03-074-dihedral-class}{Dihedral two-groups have every class}
\begin{proof}
\emph{Group structure of \(G_m\).} Multiplication by \(\pm1\) is an
additive automorphism of \(\mathbb Z/m\mathbb Z\), and
\(\varepsilon(\eta c)=(\varepsilon\eta)c\). Hence
\(((a,\varepsilon)(b,\eta))(c,\gamma)=(a+\varepsilon b+\varepsilon\eta c,
\varepsilon\eta\gamma)=(a,\varepsilon)((b,\eta)(c,\gamma))\). The element
\((0,1)\) is a two-sided identity, and
\((a,\varepsilon)(-\varepsilon a,\varepsilon)=(a-\varepsilon^2a,1)=(0,1)\),
\((-\varepsilon a,\varepsilon)(a,\varepsilon)=(-\varepsilon a+\varepsilon
a,1)=(0,1)\).

\emph{The case \(k=1\).} In \(\mathbb Z/2\mathbb Z\), \(-b=b\), so the
product in \(G_2\) is \((a+b,\varepsilon\eta)\), which is symmetric in its
two factors: \(G_2\) is abelian. It has four elements, so \(\lambda_1\)
fails and \(\lambda_2\) holds. By
\cref{thm:p03-w03-p03-074-class-law}, item 1, \(Z_0(G_2)\neq G_2\) and
\(Z_1(G_2)=G_2\): the class is \(1\).

\emph{The centre for \(m=2^k\), \(k\ge2\).} Here \(m\ge4\). The element
\((a,\varepsilon)\) is central iff
\((a+\varepsilon b,\varepsilon\eta)=(b+\eta a,\eta\varepsilon)\), that is
\(a+\varepsilon b=b+\eta a\) in \(\mathbb Z/m\mathbb Z\), for all
\((b,\eta)\). Taking \(b=0\), \(\eta=-1\) gives \(2a=0\), so
\(a\in\{0,m/2\}\). Taking \(b=1\), \(\eta=1\) gives \(\varepsilon=1\) in
\(\mathbb Z/m\mathbb Z\); since \(m\ge4\) does not divide \(2\),
\(-1\neq1\) there, so \(\varepsilon=1\). Conversely, for
\(a\in\{0,m/2\}\) and \(\varepsilon=1\) the condition reads \(a=\eta a\),
which holds for \(\eta=1\) and, as \(2a=0\), for \(\eta=-1\).

\emph{The reduction map.} Since \(m/2\) divides \(m\), reduction
\(\bar{\ }:\mathbb Z/m\mathbb Z\to\mathbb Z/(m/2)\mathbb Z\) is a
well-defined surjective ring homomorphism, and
\(\rho(a,\varepsilon)=(\bar a,\varepsilon)\) satisfies
\(\rho((a,\varepsilon)(b,\eta))=(\bar a+\varepsilon\bar b,\varepsilon\eta)
=\rho(a,\varepsilon)\rho(b,\eta)\). It is surjective, and
\(\rho(a,\varepsilon)=(0,1)\) iff \(\varepsilon=1\) and
\(a\in\{0,m/2\}\), so \(\ker\rho=Z(G_m)\).

\emph{The class.} By \cref{thm:p03-w03-p03-074-class-law}, items 1 and 2,
a group has class exactly \(c\) iff, for every \(j\ge1\),
\(\lambda_j\) holds exactly when \(j\ge c+1\). We prove this for
\(G_{2^k}\) with \(c=k\) by induction on \(k\); the case \(k=1\) is
done. Let \(k\ge2\). By \cref{thm:p03-w03-p03-074-centre-step}, item 2,
applied to \(\rho:G_{2^k}\to G_{2^{k-1}}\), for \(j\ge1\),
\(G_{2^k}\models\lambda_{j+1}\) iff \(G_{2^{k-1}}\models\lambda_j\), which
by the induction hypothesis holds iff \(j\ge k\), that is
\(j+1\ge k+1\). Also \(\lambda_1\) fails in \(G_{2^k}\), which is
nontrivial. So \(\lambda_j\) holds in \(G_{2^k}\) exactly when
\(j\ge k+1\), and the class is \(k\).
\end{proof}
\end{annexitem}

\begin{annexitem}{annex:p03-w03-p03-074-not-elementary}{thm:p03-w03-p03-074-not-elementary}{Nilpotency is not first-order}
\begin{proof}
Write \(P\) for the ultraproduct. Each group axiom holds in every factor,
the index set belongs to \(\mathcal U\), and \cref{thm:los} shows
that \(P\) is a group. Fix \(c\ge0\). For \(k>c\) the group
\(G_{2^k}\) has class \(k>c\) by
\cref{thm:p03-w03-p03-074-dihedral-class}, so by
\cref{thm:p03-w03-p03-074-class-law}, item 1, it satisfies
\(\neg\lambda_{c+1}\). The set \(\{k:k>c\}\) is cofinite, hence belongs to
the nonprincipal ultrafilter \(\mathcal U\) by
\cref{lem:nonprincipal-ultrafilter}, item 3, and \cref{thm:los}
gives \(P\models\neg\lambda_{c+1}\). By
\cref{thm:p03-w03-p03-074-class-law}, item 1, \(Z_c(P)\neq P\). As
\(c\) was arbitrary, \(P\) is not nilpotent.

Suppose a set \(\Sigma\) of sentences had exactly the nilpotent groups as
models. Each \(G_{2^k}\) is nilpotent, so every \(\sigma\in\Sigma\) holds
in every factor, and \cref{thm:los} gives \(P\models\Sigma\). Then
\(P\) would be nilpotent, which was just refuted.
\end{proof}
\end{annexitem}
